\documentclass[11pt]{article}
\usepackage[margin=1in]{geometry}
\usepackage{amsmath,amssymb,amsthm,hyperref}
\hypersetup{hidelinks}
\setlength{\emergencystretch}{2em}
\newtheorem{proposition}{Proposition}
\title{Disproof of Conjecture 00000002164\\Positive 2-factor counts cannot decay to zero}
\author{gaochengzhecpu}\date{}
\begin{document}\maketitle
\noindent\textbf{Claim addressed.} The source concerns the 2-factor count of cubic graphs and proposes $(c/\pi)^N$, with $c=4^{1/3}$, as its asymptotic expression. Under the literal interpretation as an asymptotic for the unnormalized number of 2-factors, this is impossible: the proposed base is less than one, whereas an explicit infinite family of connected cubic graphs has at least one 2-factor at every size in the family.

A 2-factor is a spanning subgraph in which every vertex has degree two. For $m\ge3$, let
\[
 G_m=C_m\mathbin{\square}K_2
\]
be the prism graph. Its vertices are $(i,b)$, where $i\in\mathbb Z/m\mathbb Z$ and $b\in\{0,1\}$. There are cycle edges within each layer and one edge from $(i,0)$ to $(i,1)$ for each $i$. Let $F_m$ denote its number of 2-factors.

\begin{proposition}
The graph $G_m$ is a connected simple cubic graph on $2m$ vertices, and $F_m\ge1$. With $b=4^{1/3}/\pi$, the sequence $F_m$ is not $O(b^{2m})$ as $m\to\infty$.
\end{proposition}
\begin{proof}
Each vertex has two distinct neighbors in its cycle layer and one in the other layer, hence degree three. The cycles and the edges joining their corresponding vertices make the graph connected. Deleting all edges between the layers leaves the two disjoint cycle layers, whose union spans $G_m$. This union is a 2-factor, proving $F_m\ge1$.

The proposed base satisfies
\[
 0<b=\frac{4^{1/3}}{\pi}
 \le\frac{\sqrt4}{\pi}=\frac2\pi<1.
\]
Here monotonicity of real powers of $4>1$ gives the middle inequality, and $\pi>3>2$ gives the last one. Consequently $b^{2m}\to0$. If $F_m=O(b^{2m})$, there would be a constant $C>0$ with
\[
 1\le F_m\le Cb^{2m}
\]
for all sufficiently large $m$. The right side tends to zero, a contradiction.
\end{proof}

In particular the displayed expression cannot be an asymptotic equivalent of these counts, nor can multiplying it by a fixed constant make it one. The failure is asymptotic along an infinite connected family, not just at a small graph.

\noindent\textbf{Scope.} The source says ``2-factor count'' and supplies no normalization, averaging law, or additional restriction on cubic graphs. The counterexample addresses that unnormalized, unrestricted counting assertion. No assertion is made about a differently defined probability, normalized count, or ensemble average. It is not necessary to claim that every cubic graph has a 2-factor: only the displayed family is used. The argument does not evaluate the Fisher--Kasteleyn mapping; any count-preserving mapping must still respect the positive integer lower bound.

\section*{Lean correspondence and reproduction}
The formal graph \texttt{prism n} is the actual Mathlib Cartesian graph product of \texttt{cycleGraph (n+3)} and the complete graph on \texttt{Fin 2}. Theorems prove its connectedness and degree three. The subgraph \texttt{rings n} is the product with the empty graph on the same two vertices. The proof establishes its inclusion in the prism and degree two at every vertex.

The type \texttt{TwoFactor n} consists of \emph{all} simple graphs on the prism's vertex type that are subgraphs of the prism and have degree two everywhere. Its natural cardinality is the actual count. Finiteness and an explicit inhabitant give \texttt{count n} $\ge1$. The actual vertex cardinality is proved equal to $2(n+3)$; thus the exponent is the number of vertices, not a substituted size parameter.

The formal base uses \texttt{Real.rpow} and \texttt{Real.pi}. The proof establishes convergence of the proposed expression to zero and failure of convergence to zero for the counts. The final theorem negates the actual \texttt{Asymptotics.IsBigO} statement at \texttt{atTop}, using the exponent \texttt{vertexCount n}. No asymptotic value or graph count is assumed by definition.

Inside \texttt{lean/}, run \texttt{lake build}, followed by
\begin{center}\texttt{lake env lean -DwarningAsError=true Main.lean}.\end{center}
Lean is pinned to 4.19.0 and Mathlib to
\begin{center}\small\texttt{c44e0c8ee63ca166450922a373c7409c5d26b00b}.\end{center}
All transitive dependencies are pinned. No auxiliary numerical computation is needed. Build logs, axiom dependencies, source hashes and PDF checks are in \texttt{verification/}. Author and delegated-agent review records are separate; no external independent review is claimed.

\begin{thebibliography}{9}
\bibitem{source} The Last Math Competition,
\href{https://github.com/The-Last-Math-Competition/The-Last-Math-Competition/blob/main/conjectures/00000002164.md}{Conjecture 00000002164 (original bilingual source)}.
\end{thebibliography}
\end{document}
